\section{Chương~\ref{sec:dofaalt}. Các lý thuyết khác về \nt{DOPA}}

Mỗi bức tranh mở rộng dựa trên những phép đo dopamine trực tiếp của riêng nó (độ phân tán của các điểm đảo chiều, đáp ứng với điều khó chịu và điều mới, sự tăng chậm, đáp ứng khi đổi vị), nhưng các công thức giải thích chúng là lý thuyết, và giữa hai trong số đó các thí nghiệm hiện vẫn mâu thuẫn nhau.

{\footnotesize
\setlength{\tabcolsep}{3pt}\renewcommand{\arraystretch}{1.15}
\begin{longtable}{>{\raggedright}p{0.5cm}>{\raggedright}p{4.0cm}>{\raggedright}p{5.6cm}>{\raggedright}p{2.3cm}>{\raggedright\arraybackslash}p{1.7cm}}
\toprule
TT & Khẳng định & Thí nghiệm và điều đã đo & Nguồn & Trạng thái \\
\midrule\endfirsthead
\toprule
TT & Khẳng định & Thí nghiệm và điều đã đo & Nguồn & Trạng thái \\
\midrule\endhead
1 & Quy tắc bất đối xứng~\eqref{eq:asymmetric} hội tụ về điểm~\eqref{eq:expectile}; khi $\kappa=1/2$ đó là trung bình, với một giọt có xác suất $1/2$ thì $V=\kappa$ (hình~\ref{fig:expectiles}) & Điểm~\eqref{eq:expectile} là expectile, nghiệm của bài toán bình phương tối thiểu với trọng số khác nhau cho độ lệch dương và âm; với ví dụ trong chương, phép suy ra có ngay trong chương. & \cite{NeweyPowell1987}; suy ra trong chương & lý thuyết \\
2 & Các nơ-ron \nt{DOPA} khác nhau có điểm đảo chiều khác nhau, và chúng được sắp theo độ bất đối xứng (mệnh đề~\ref{prop:expectile}) & Chuột nhắt, nơ-ron \nt{DOPA} vùng mái bụng được đánh dấu quang di truyền, phần thưởng có kích thước khác nhau: điểm mà đỉnh vọt chuyển thành hõm sụt khác nhau ở các nơ-ron; ở nơ-ron có đáp ứng bất đối xứng hơn thì điểm này cao hơn; từ đáp ứng của nhóm có thể khôi phục hình dạng phân bố phần thưởng. Việc đây là chức năng chính của sự phân tán là giả thuyết. & \cite{Dabney2020} & đã đo \\
3 & Một phần nơ-ron \nt{DOPA} đáp ứng với điều khó chịu bằng đỉnh vọt & Khỉ: một số nơ-ron \nt{DOPA} bị kích thích bởi tín hiệu báo trước phần thưởng và bị ức chế bởi tín hiệu báo trước luồng khí thổi vào mắt, số khác bị kích thích bởi cả hai; các nhóm nằm ở những phần khác nhau của não giữa. & \cite{Matsumoto2009, BrombergMartin2010} & đã đo \\
4 & Độ lớn của sai số hoặc độ mới đặt tốc độ học & Các công trình được trích không có thí nghiệm trực tiếp nào trong đó tín hiệu tầm quan trọng của nơ-ron \nt{DOPA} làm thay đổi tốc độ học; bài tổng quan đề xuất vai trò tín hiệu ``chú ý, quan trọng''. & \cite{BrombergMartin2010} & giả thuyết \\
5 & Một dây riêng cho trục nguy hiểm & Chuột nhắt: nơ-ron \nt{DOPA} đi tới phần đuôi của thể vân đáp ứng với kích thích mới và mạnh chứ không với phần thưởng; phá hủy chúng làm yếu đi hành vi tránh một vật mới. & \cite{Menegas2018} & đã đo \\
6 & Thời gian hành động tối ưu $\tau_a^*=\sqrt{C/\bar R}$~\eqref{eq:vigor} & Cực tiểu của tổng $C/\tau_a+\bar R\tau_a$; suy ra trong chương. Trong mô hình gốc, tốc độ hành động do giá của thời gian quyết định, bằng phần thưởng trung bình. & \cite{Niv2007}; suy ra trong chương & lý thuyết \\
7 & Mức dopamine chậm mang $\bar R$ và đặt tốc độ hành động (mệnh đề~\ref{prop:multiplex}) & Chuột cống, vi thẩm tách và đo vôn-ampe ở nhân accumbens: mức dopamine theo từng phút bám theo tần suất phần thưởng và tốc độ hành động. Ở người, L-DOPA làm tăng sự phụ thuộc của thời gian phản ứng vào tần suất phần thưởng trung bình. Việc mức chậm đúng bằng $\bar R$ chưa được chứng minh. & \cite{Hamid2016, Beierholm2013vigor} & giả thuyết \\
8 & Mức chậm được tạo từ hai nguồn: lượng phóng thay đổi ở đầu ra khi tần số xung không đổi, nhưng sự tăng chậm cũng có trong chính tần số & Chuột cống, cùng một nhiệm vụ: lượng phóng ở nhân accumbens bám theo kỳ vọng phần thưởng thay đổi, còn tần số xung của các nơ-ron \nt{DOPA} vùng mái bụng đã được nhận dạng thì không. Chuột nhắt trong hành lang ảo (dòng~10): sự tăng dần đều khi tiến gần phần thưởng có cả trong xung của các nơ-ron \nt{DOPA} đã được nhận dạng. & \cite{Mohebi2019, Kim2020} & đã đo \\
9 & Dopamine tăng dần đều trong khi con vật tiến gần một phần thưởng ở xa & Chuột cống trong mê lộ, đo vôn-ampe ở thể vân: nồng độ dopamine tăng trong nhiều giây khi tiến gần đích, độ dốc phụ thuộc vào khoảng cách và kích thước phần thưởng. & \cite{Howe2013ramp, Hamid2016} & đã đo \\
10 & Sự tăng là sai số $\delta$ khi ước lượng chính xác dần~\eqref{eq:ramp}; bước nhảy trong hành lang cho đỉnh vọt (hình~\ref{fig:ramp}) & Công thức và con số $\delta=0.75\,V$ mỗi giây được suy ra trong chương. Thí nghiệm: chuột nhắt trong hành lang ảo, được dời tức thời lại gần đích; xung ở thân tế bào, calci ở thân và ở đầu ra, và nồng độ trong thể vân đáp ứng như sai số chứ không như kỳ vọng $V$. Cách giải thích nào trong hai cách là đúng ở mọi đầu ra thì vẫn đang được bàn cãi. & \cite{Kim2020} & đã đo \\
11 & Nhìn lại phía sau: dopamine báo thước đo liên hệ nhân quả~\eqref{eq:retro} & Chuột nhắt, lượng dopamine phóng ở nhân accumbens trong những thí nghiệm mà lý thuyết này và sai số~\eqref{eq:ctd} dự đoán khác nhau: trong một số thí nghiệm, lượng phóng tuân theo lý thuyết thứ nhất. & \cite{Jeong2022} & giả thuyết \\
12 & Trong quá trình học, đáp ứng dopamine dịch dần từ phần thưởng sang tín hiệu báo trước & Chuột nhắt, tín hiệu ở đầu ra của nơ-ron \nt{DOPA} trong thể vân bụng theo tiến trình học: đỉnh vọt dần chuyển từ thời điểm có phần thưởng sang thời điểm có tín hiệu báo trước, đúng như việc học theo~\eqref{eq:ctd} đòi hỏi. & \cite{Amo2022} & đã đo \\
13 & Nơ-ron \nt{DOPA} đáp ứng khi vị của phần thưởng đổi mà giá trị vẫn như cũ & Chuột cống, ghi nơ-ron vùng mái bụng: thay nước ép bằng một loại khác có giá trị tương đương gây ra đáp ứng, dù sai số~\eqref{eq:ctd} bằng không. & \cite{Takahashi2017} & đã đo \\
14 & Đỉnh vọt nhân tạo dạy mối liên hệ giữa hai kích thích trung tính & Chuột cống, thí nghiệm ghép cặp trước hai kích thích không có phần thưởng: kích hoạt quang di truyền nơ-ron \nt{DOPA} lúc chuyển từ kích thích thứ nhất sang thứ hai tạo ra mối liên hệ, tắt chúng thì cản trở việc học liên hệ đó. & \cite{Sharpe2017} & đã đo \\
15 & Vectơ sai số theo đặc trưng~\eqref{eq:vectortd} & Lý thuyết về sai số dự đoán đặc trưng; chưa có kiểm tra trực tiếp dạng vectơ. & \cite{Gardner2018} & giả thuyết \\
16 & Trong cùng một nhiệm vụ, các nơ-ron \nt{DOPA} khác nhau mang những đại lượng khác nhau & Chuột nhắt trong hành lang ảo, chụp hai photon nơ-ron \nt{DOPA}: một số liên quan tới phần thưởng, số khác tới tốc độ và hướng rẽ, số khác nữa tới tín hiệu báo trước và độ chính xác của lựa chọn. & \cite{Engelhard2019} & đã đo \\
17 & Bó dây thay cho một dây (mệnh đề~\ref{prop:bundle}) & Tổng hợp các dòng~2--16: mỗi cách phân tách đã được đo riêng; chưa có thí nghiệm cho bức tranh thống nhất rằng tất cả đều là phần của một tín hiệu. & --- & giả thuyết \\
\bottomrule
\end{longtable}}
